\frontchapter{Au lecteur}
\origpage{8}{前言 004}
\begingroup\itshape
\begin{center}\bfseries « Non abbiate paura »\end{center}
Français de nationalité, chinois de cœur, voilà de nombreuses années que j’enseigne les sciences du langage dans de différentes universités chinoises (北京语言大学，中国人民大学 et 北京外国语大学) à un public d’étudiants de niveau licence, master et doctorat. Au cours de ces années de professorat, j’ai retenu un enseignement : la linguistique est la terreur absolue des étudiants ! Oui, la linguistique fait peur. Mais ce qui est intéressant, c’est que cette peur est en grande partie irrationnelle, fondée sur la mauvaise réputation qui précède cette matière souvent jugée difficile. Pour preuve, chaque début de semestre, avant de commencer le premier cours, en préambule, je fais le même petit test auprès de mes étudiants : je leur demande de qualifier la linguistique. Après quelques secondes d’hésitation, -- ils redoutent de me faire de la peine, les langues se délient. Le tiercé gagnant est invariablement : difficile, ennuyeux, et inutile. Rien de très réjouissant, me direz-vous. Or, loin de me chagriner, ces réponses m’amusent beaucoup car ces adjectifs reposent sur un a priori plus que sur un constat a posteriori. En revanche, ce qui est révélateur, c’est qu’à la moitié du semestre, les sourires et la bonne humeur remplaceront le silence et les mines blêmes du premier cours. Tandis que le taux de 100 \% de réussite à l’examen est véritablement ardu, c’est une preuve évidente de l’engouement et de l’adhésion suscités par cette matière à la triste réputation.

Alors, cher lecteur, avant de faire appel à votre raison pour vous démontrer que cette triple réputation infamante n’est pas justifiée, je vais d’abord faire appel à vos émotions et tenter de vous rassurer. Pour cela, et bien que je ne sois pas le pape des sciences du langage, j’ai emprunté au pape Jean Paul II son célèbre « non abbiate paura », ce qui, en italien, signifie « n’ayez pas peur ».

\begin{center}\bfseries Voici pourquoi\end{center}
Au moment d’acheter le livre que vous tenez dans vos mains, vous vous êtes peut-être demandé pourquoi étudier la linguistique. Avant de vous répondre, je voudrais vous poser une question : si vous avez un animal de compagnie (chien, chat, tortue, poisson rouge), demandez-vous ce qui vous différencie de lui. La réponse, bien plus simple que vous l’imaginez, puisque nous sommes nous-mêmes des animaux, tient en trois mots : art, pensée, langage. C’est grâce à l’art (c’est-à-dire la création désintéressée du beau), à la conscience de soi (votre chat ne sait pas qu’il est un chat), et le langage que nous dépassons notre animalité et acquérons notre humanité. C’est dire l’importance du langage.

Si, pour utiliser un mot compliqué, le langage est consubstantiel à l’espèce humaine, la langue que nous parlons, quant à elle, nous définit dans une grande mesure. Pour le poète et écrivain roumain d’expression française Emil Cioran (1911-1995), on n’habite pas un pays, on habite une langue. Le philosophe français Jacques Derrida (1930-2004) a exprimé une idée semblable : « Je n’ai qu’une langue [...] Je l’appelle ma \origcont{9}{前言 005}demeure, je le ressens comme tel, j’y reste et je l’habite ».

Notre langue est donc ce qu’il y a de plus proche de nous. Or, paradoxalement, du fait même de cette proximité, on ne s’y intéresse pas, on n’y pense même pas. On parle comme on respire, sans s’en rendre compte. Le linguiste et anthropologue américain Edward Sapir (1884-1939), que nous rencontrerons dans le chapitre 14, ne dit pas autre chose : « Le langage est quelque chose de tellement familier dans notre vie quotidienne que l’on s’arrête rarement pour essayer de le définir. Il est aussi naturel pour l’homme de parler que de marcher ou de respirer ». Étudier les sciences du langage permet donc d’introduire la conscience et la réflexion dans une faculté presque inconsciente et à laquelle on ne songe jamais. Voilà donc une première raison de vous plonger dans la lecture de ce livre.

La deuxième raison vous concerne directement, cher lecteur. Durant vos longues années d’apprentissage du français, vous avez appris les quatre savoirs-faire\ednote{原文作 « savoirs-faire »。此复合名词应写作 « savoir-faire »，复数形式不变，故此处为 « les quatre savoir-faire »。} fondamentaux : écouter, parler, lire et écrire. Le cours de linguistique a pour but, en s’appuyant sur votre pratique des langues, de vous amener à réfléchir sur le langage, et donc à prendre à la fois du recul et de la hauteur. Cette combinaison idéale de pratique et de réflexion viendra parachever votre cursus universitaire.

Pour Blaise Pascal (1623-1652)\ednote{原文生卒年如此。Pascal 的卒年应为 1662，故应作 « 1623-1662 »；参见法国国家图书馆作者资料。\ecite{1}}, il y a deux sortes d’esprit : l’esprit de finesse et l’esprit de géométrie, c’est-à-dire la raison poétique et la raison scientifique. Si vos cours de littérature française ont développé votre esprit de finesse, le cours de linguistique s’adresse, lui, à votre esprit de géométrie. J’ajoute que la culture scientifique fait partie de la culture. Le mathématicien français Cédric Villani, lauréat en 2010 de la médaille Fields (équivalent du prix Nobel pour les mathématiques) recommande d’intégrer les sciences dans la culture, et je souscris à son point de vue. Contrairement à une idée trop répandue, la culture n’est pas que littéraire, et ce livre va considérablement développer votre culture scientifique et vous conduire vers de nouveaux territoires demeurés jusqu’alors insoupçonnés.

Les raisons d’étudier la linguistique, aussi valables soient-elles, ne peuvent occulter la difficulté de son apprentissage. N’est-ce pas le cas de tout apprentissage ? Identifier les raisons de cette difficulté, c’est déjà l’avoir à moitié surmontée, alors commençons par les raisons intrinsèques.

En français, vous l’avez noté dans le titre de notre livre, la linguistique est appelée « sciences du langage ». Cela nous donne trois pistes de réflexion :

Tout d’abord, on parle de « science », donc de rigueur, de méthode, d’objectivité, de logique. Cela peut faire peur au lecteur non averti, mais au fond c’est un avantage. Dans ce livre, pas de « blah blah » : rien que des concepts clairs et bien définis, donc aisés à comprendre.

Ensuite, le « s » accolé au mot « science » nous indique que la linguistique, un peu comme une poupée russe, est en fait un ensemble de sciences. En général, on retient six domaines (qui constituent le cœur et la deuxième partie de notre livre) : la phonétique, la phonologie, la morphologie, la syntaxe, la sémantique, et la pragmatique. Il y a bien d’autres\ednote{原文缺少代词 « en »；应作 « Il y en a bien d’autres »，意为还存在其他领域。}, et notre livre comporte 16 chapitres. Ici encore, cela peut sembler compliqué. Mais non en fait ! La linguistique n’est pas compliquée mais complexe. La langue est un système de systèmes : ces six domaines sont imbriqués les uns dans les autres, et, comme l’a dit Blaise Pascal, qui en plus d’être philosophe était un scientifique de génie, « je ne peux comprendre un tout que si je connais \origcont{10}{前言 006}particulièrement les parties, mais je ne peux comprendre les parties que si je connais le tout ». Cela s’applique tout à fait au langage, et la complexité des sciences du langage, loin d’être un défi, en facilite grandement leur apprentissage. Un plat de spaghetti, c’est compliqué. Un circuit électronique est complexe. Pensez aux sciences du langage comme à un circuit, et vous comprendrez ce que je veux dire.

Comme leur nom l’indique, les sciences du langage portent sur le langage, qui, contrairement aux langues, est une faculté universelle et abstraite. Cela, je le reconnais, peut constituer une réelle difficulté car il vous faudra faire un effort d’abstraction pour vous dégager des langues particulières et remonter vers des principes généraux, voire universaux. Un facteur aggravant, si je puis dire, est que au sein de nos facultés et départements de français, la très grande majorité des étudiants ont eu un parcours littéraire et sont imperméables à la pensée logique et abstraite. Conscient de cette difficulté, et pour vous permettre d’appréhender ces nouveaux territoires de pensée, le chapitre 4 est entièrement consacré à la logique et à la philosophie analytique. Sa lecture vous sera très profitable pour aborder les autres chapitres.

Une autre difficulté des sciences du langage réside dans l’utilisation d’un jargon qui, je le reconnais bien volontiers, peut faire peur de prime abord. Les linguistes justifient l’utilisation de ce métalangage (lui-même un bel exemple de jargon) par le fait que la langue de tous les jours, celle que nous utilisons vous et moi, n’est pas appropriée, pas suffisamment précise, trop équivoque pour être utilisée dans une discipline scientifique faite, on l’a vu, de rigueur et de précision. Alors, oui, c’est vrai : vous allez rencontrer dans ce livre un nombre assez impressionnant de termes que même les Français les plus érudits ne connaissent pas : glossolalie, allophone, hyperpolyglossie, archisémème, isotopie, aphérèse, prosthèse, parataxe, exophore, amphibologie, métaplasme, etc. : la liste est très longue. Mon conseil est de ne pas vous laisser impressionner par les mots : l’important est de comprendre les concepts auxquels ils renvoient. Pour vous aider à les démystifier, je les ai tous sans exception « traduits » dans notre langue de tous les jours. Ensuite, pour la très grande majorité d’entre eux, j’ai indiqué leur étymologie grecque et/ou latine, de façon à vous permettre d’en deviner le sens. Ici encore, la difficulté n’est qu’apparente.

Parce que le langage concerne tous les domaines du savoir, la linguistique est une discipline éminemment transversale. Cela veut dire que notre livre ne va pas se limiter aux six domaines des sciences de langage évoqués plus haut. Nous allons devoir faire appel à un grand nombre d’autres sciences et disciplines : histoire, mythologie, religion, philosophie, logique, ethnologie, biologie, psychologie, anthropologie, paléontologie, éthologie, biologie, neurologie, génétique, sémiologie, cybernétique, informatique, communication, intelligence artificielle, physique, chimie, et même, dans notre dernier chapitre, musique ! Cet « inventaire à la Prévert » a de quoi inquiéter, mais, encore une fois, il faut le voir comme facteur positif, car on peut tout reprocher à la linguistique, sauf d’être ennuyeuse. Puisqu’on a évoqué la musique, le compositeur italien Gioacchino Rossini, (1792-1868) pensait que « tous les genres sont bons hors le genre ennuyeux ». Il a cent fois raison, et je puis vous garantir que le livre que vous avez entre vos mains est divertissant.

J’espère avoir démontré que la difficulté est plus une croyance tenace qu’un fait avéré. Mais il n’en demeure pas moins que tout étudiant sera confronté à un triple défi : devoir assimiler de très nombreux concepts, appréhender des domaines très divers, dans un laps de temps assez bref (en général, un à deux semestres) ! Néanmoins, là encore, je vais vous rassurer : en qualité d’étudiants en langues étrangères, vous êtes déjà \origcont{11}{前言 007}des linguistes. Seulement, vous ne le savez pas. Vous êtes un peu comme Monsieur Jourdain du Bourgeois Gentilhomme de Molière : « Il y a plus de quarante ans que je dis de la prose sans que j’en susse rien ! » Vous aussi, cher lecteur, il y a des années que vous êtes linguistes sans que vous en sachiez rien.

Autre facteur qui vous est favorable, et non des moindres, vous êtes aussi polyglottes : vous parlez le mandarin, l’anglais, le français, et il est fort possible que vous parliez un dialecte du chinois, et probable que vous connaissiez des éléments d’une quatrième langue apprise, en option ou par plaisir dans votre temps libre. Ceci va grandement vous aider dans votre apprentissage des sciences du langage : s’il n’est pas à proprement parler nécessaire de connaître plusieurs langues pour faire de la linguistique, le fait d’en connaître plusieurs, de pouvoir les comparer, de voir ce qui les rapproche malgré leurs différences apparentes, est fondamental dans la réflexion linguistique. Dans le chapitre 5, consacré justement à la linguistique comparative, vous verrez que les linguistes comparatistes étaient tous des « hyperpolyglottes » capables de parler plusieurs dizaines de langues, voire plus d’une centaine pour certains. Cela vous laisse rêveur ? Moi aussi !

Si je n’ai pas réussi à vous convaincre, alors il ne me reste plus qu’à recourir à un argument massue et faire appel au philosophe romain Sénèque : « Ce n’est pas parce que les choses sont difficiles que nous n’osons pas les faire, c’est parce que nous n’osons pas les faire qu’elles sont difficiles. »

Qu’en est-il à présent de l’inutilité, deuxième critique fréquemment adressée à la linguistique ? Avant d’y répondre, je note qu’on pourrait adresser la même critique à la littérature, l’histoire, la philosophie, la théologie, la sociologie, l’art, et à toutes les sciences et disciplines non appliquées, et la liste est potentiellement longue. Ensuite, si par inutile on entend qu’elle ne permet pas de trouver un travail, je répondrai que l’université n’a pas pour vocation première de vous aider à trouver un travail, mais de former des esprits, des têtes bien faites et bien pleines. Et d’ailleurs, dans le chapitre 3, vous verrez que dans les universités françaises au Moyen Âge les « arts libéraux », et en particulier la grammaire, la rhétorique, et la dialectique, regroupées sous le nom de « trivium », étaient les plus prisés précisément parce que visant une connaissance désintéressée, et par conséquence, supérieure. À l’opposé, les arts « serviles », c’est-à-dire utiles, étaient déconsidérés. Même si très récemment de nouveaux métiers sont apparus dans les domaines de l’intelligence et des langages artificiels, du traitement automatique du langage, de la linguistique informatique, ou encore de la traduction automatique, je ne pense pas qu’il faille raisonner en termes d’utilité. Il y a plus de 400 ans, le poète oublié François de Malherbe (1555-1628), bien que poète officiel du roi Henri IV, était pleinement conscient de son inutilité : « Un bon poète n’est pas plus utile à l’État qu’un bon joueur de quille ». Dans le même esprit, je veux bien que l’on dise qu’un bon linguiste n’est pas plus utile à l’État qu’un bon joueur de quille. Car comme l’a très bien dit un autre poète, le « poète impeccable » Théophile Gautier (1811-1872) : « Tout ce qui est utile est laid. Il n’y a vraiment de beau que ce qui ne peut servir à rien. » Vous pourrez me rétorquer que, si la poésie et les sciences du langage sont aussi inutiles l’une que l’autre, il existe une différence de taille : la poésie et l’art en général sont des vecteurs d’émotion, alors que la linguistique a la froideur de la science, la distance des concepts et absence de chair de l’abstraction. C’est tout à fait vrai, et c’est là qu’intervient le rôle du professeur.

\origpage{12}{前言 008}
\begin{center}\bfseries Voilà comment\end{center}
La linguistique est une science froide et éloignée de vos préoccupations. Comment ne pas le comprendre ? Quand on est un étudiant d’une vingtaine d’années, difficile de se passionner pour la « sémiotique phanéroscopique » de Charles Sanders Peirce (chapitre 6) ou la « grammaire générative transformationnelle » de Noam Chomsky (chapitre 14). C’est bien naturel ! Pour le professeur de linguistique, le modus operandi tient donc en deux mots : rapprocher et humaniser.

D’abord, rapprocher nos étudiants des textes originaux. Pour cela, notre livre contient un très grand nombre d’extraits d’œuvres originales en français, puisque la France a la particularité d’avoir produit un grand nombre de linguistes de premier plan. Pour les étudiants de licence, master et a fortiori de doctorat, le niveau de langue n’est plus un problème : ils peuvent donc aborder directement les œuvres originales, ce qui, je le constate chaque semestre, est très gratifiant pour eux. Les deux chapitres liminaires de notre livre sont l’occasion de faire « lire dans le texte » à nos étudiants de longs extraits des œuvres de Montaigne, Rousseau, Descartes, et Bergson portant sur le langage. Viendront ensuite les linguistes (Émile Benveniste, André Martinet, Ferdinand de Saussure, Henri Frei, Charles Bally, Lucien Tesnière, etc.).

Il faut ensuite rapprocher nos étudiants de la linguistique contemporaine. Déjà, au 17\textsuperscript{e} siècle, le moraliste français Jean de la Bruyère (1645-1696) écrivait : « Tout est dit, et l’on vient trop tard depuis plus de sept mille ans qu’il y a des hommes qui pensent. » C’est encore plus vrai plus de trois cent cinquante ans après. Si les quatre chapitres de la première partie présentent la genèse de la linguistique, depuis le mythe de Babel jusqu’à la naissance de la linguistique moderne avec Ferdinand de Saussure, les quatre chapitres de la troisième partie sont quant à eux consacrés aux développements les plus récents : sociolinguistique, grammaire générative, biolinguistique, linguistique cognitive, etc. Notre livre s’achève avec un entretien de décembre 2017 avec le neurobiologiste Antonio Damasio : difficile de faire plus actuel ! L’idée est de montrer à nos étudiants que la linguistique est une science très jeune, en plein développement, et fourmillante de concepts nouveaux. Avec infiniment de respect, je vais corriger un peu Jean de la Bruyère, et dire : « En linguistique, tout n’a pas été dit. » Nos étudiants chercheurs ont encore de nombreux pans de la linguistique à défricher.

Humaniser aussi et encore. Il s’agit ainsi de donner de la chaleur à une science froide. L’idée est simple : avant de présenter les concepts, présenter les hommes. Les logiciens, anthropologues, linguistes, philosophes, sémiologues, etc. présentés dans notre livre sont tous des savants polymathes dont le savoir encyclopédique et le parcours intellectuel sont très inspirants. Dans notre livre, les concepts et théories sont systématiquement précédés d’une présentation de leurs auteurs. C’est nécessaire pour donner un peu de « chair » et de vie, mais ce n’est pas suffisant : donner à voir et entendre, c’est beaucoup mieux. Et nous avons de la chance : à notre époque moderne, il suffit d’un clic sur Internet pour avoir accès à des cours, entretiens et conférences sous format vidéo donnés par les plus illustres linguistes et spécialistes du langage modernes et contemporains. Quoi de mieux que de les écouter nous exposer leurs théories ?

Voilà, cher lecteur, ce que je souhaitais vous dire avant que vous vous lanciez dans la lecture de cet ouvrage. Dans son Gai Savoir (1882) le philosophe et poète allemand Nietzsche a écrit : « N’est-ce pas une chose extrêmement plaisante que de voir les philosophes les plus sérieux, si sévères qu’ils soient le reste du temps avec toute certitude, en appeler sans cesse à des sentences de poètes pour assurer force et crédibilité à leur \origcont{13}{前言 009}pensée ? » Certes, je ne suis pas philosophe, juste linguiste, mais si j’ai moi-aussi utilisé les sentences des poètes pour assurer force et crédibilité à mon argument, c’était dans le seul but de vous rassurer. J’espère avoir réussi. Je souhaite surtout que ce livre vous donne les clefs pour aborder avec curiosité et enthousiasme le « gai savoir » de la linguistique.

Bonne lecture.
\begin{flushright}Lionel Spinosa\end{flushright}
\endgroup
